\chapter{Anhang}
\label{ch:Anhang}

\img{plots/Positionen_der_Minima2a}{Abstände der Extrema in Abhänigkeit ihrer Ordnungszahl. Maxima sind falsch eingezeichet; sind um 1 zu klein.}
Leider ist die Ordnungszahl der Maxima falsch bestimmt, diese ist chronisch um $1$ zu klein. Dies ist insbesondere der \ctab{vergleich_ordnungszahl_theo_ex_einzelspalt} zu entnehmen und auch logisch zu begründen. Denn die Iteration berücksichtigt nicht, dass das Hauptmaximum übersprungen werden muss.
\begin{table}[h!]
    \caption{Vergleich der experimentell bestimmten Ordnungszahlen un der theorisch erwatreten. Zudem wurden die experimentellen Daten korrigiert, den diese warern alle um 1 zu klein. Die Fehlerherkunft konnte nicht ausfindig gemacht werden.}
    \label{tab:vergleich_ordnungszahl_theo_ex_einzelspalt}
    \centering
    \begin{tabular}{L C | C || C | C}
        \toprule
        n_\text{theo} & n_\text{exp} & \text{Std. Abw. } \sigma & n_\text{exp,korr} & \text{Std. Abw., korr} \sigma \\
        \midrule
        1.43 & 0.58 \pm 0.15 & 5.76 & 1.58 \pm 0.15 & 0.99 \\
        2.46 & 1.43 \pm 0.15 & 6.79 & 2.43 \pm 0.15 & 0.20 \\
        3.47 & 2.47 \pm 0.16 & 6.20 & 3.47 \pm 0.16 & 0.02 \\
        4.48 & 3.47 \pm 0.17 & 5.87 & 4.47 \pm 0.17 & 0.06 \\
        5.48 & 4.49 \pm 0.19 & 5.29 & 5.49 \pm 0.19 & 0.05 \\
        6.48 & 5.46 \pm 0.2 & 5.03 & 6.46 \pm 0.2 & 0.12 \\
        \bottomrule
    \end{tabular}
\end{table}
